\scanpage{51}
\mspage{10}
Cela fait $72=2\times36$ racines, associées aux $27\times16=6\times72$
\inserted{$=432$} couples $(\delta',\delta)$ de sommets distincts non liés

On vérifie de plus que pour \unsure{tte|les|tt} \unsure{racine|racines} $r_\alpha$, associé à un système
$\left\{\begin{aligned}\alpha&=(\alpha_1,\ldots,\alpha_6)\\\alpha'&=(\alpha'_1,\ldots,\alpha'_6)\end{aligned}\right.$, on a
\[
r_\alpha\alpha_i=(\alpha'_i-\alpha_i)\alpha_i=0-(-1)=1,\quad
\inserted{\text{de }=}\ r_\alpha(\alpha'_i)=-1
\]
\[
\left\{\begin{aligned}
r_\alpha\cdot\alpha_i&=1\\
r_\alpha\,\alpha'_j&=-1\\
r_\alpha\,\deleted{\unsure{\alpha_j}}\delta&=0\quad\text{si }\delta\in\Delta_0\text{ distinct de }\alpha_i,\alpha'_i .
\end{aligned}\right.
\]
% slanted
\margin{ceci provient du fait que $\unsure{\exists}$ \unsure{couple} \deleted{\unsure{$(\alpha_i,\alpha'_i)$|$(\alpha'_i,\alpha_i)$}} (\unsure{en} \deleted{\unsure{fait}} \unsure{il y en a ex. $4$|il y en a $4$}) tel que $\delta$ soit \unsure{lié} à $\alpha_i,\alpha'_i$ i.e. $\delta$ orthog. à $\alpha_i,\alpha'_i$ donc à leur diff. $r$}

Donc cela explicite \inserted{(pour un $\Delta$ cubique génér. ?)} comment à une racine $r$ sont associées deux parties $\alpha(r),\alpha'(r)$ de $\Delta$ : l'une des $s$ \inserted{$\in\Delta$} tels que $rs=1$ resp. $-1$ ; On trouve une corr. $1$-$1$ entre ces parties \inserted{$\alpha$} et les \deleted{parties} racines, le passage de $r$ à $-r$ revenant à échanger $\alpha$ avec le $\alpha'$ dual \ldots

Bien entendu, \deleted{les} \inserted{les} parties basiques de $\Delta$ \unsure{sont} associées à \deleted{des} \inserted{une} racine. Pour voir que l'inverse est vrai, il suffit de vérifier que les parties écrites ci-dessus \inserted{dans le cas standard} sont basiques (et ils sont groupés alors par couples de bases duales).

\deleted{Il reste \unsure{sont} \unsure{donc} que :} \deleted{\inserted{à vérifier que} dans $\widehat E(\Delta_0)$}
\[
\deleted{\alpha'_i-\alpha_i=r\quad \text{indépendant de }i, \quad\text{considérons}}
\]
\[
\begin{aligned}
&\deleted{\eta(\alpha)=\lambda-r(\alpha)=\lambda-r}\\
&\deleted{\text{et }\textstyle\sum\alpha_i=r+\zeta(\alpha)\ ;\quad\text{vérifions qu'il est \unsure{divisible|accessible} par 2 dans}}\\
&\deleted{\widehat E,\qquad\underbrace{\textstyle\sum\alpha_i}_{\zeta(\alpha)}\quad
 \text{sont }\eta(\alpha)=\tfrac12(r+\zeta(\alpha))}
\end{aligned}
\]
\deleted{On trouve dans les cas \uline{listés}}
\[
\left\{\begin{aligned}
&\deleted{\eta(r_0)=\eta(\zeta_1,\ldots\zeta_6)=\eta}\\
&\deleted{\eta(-r_0)=\eta(\zeta'_1\ldots\zeta'_6)=\eta'=5\eta-2\zeta=5\eta-2\textstyle\sum\zeta_i}\\
&\deleted{\eta(\zeta_i-\zeta_j)=\eta(\zeta_j,\zeta'_j,(\unsure{\zeta_{ih}|\xi_{ik}}))=3\eta-2\zeta_i-\textstyle\sum_{h\ne i,j}\zeta_h=(=\lambda)-\zeta_i+\zeta_j=\lambda-r.}
\end{aligned}\right.
\]
\[
\begin{aligned}
&\deleted{\eta(\alpha)}\quad\deleted{r=2\eta(\alpha)-\underbrace{\textstyle\sum\alpha_i}_{\zeta(\alpha)}\quad\text{i.e. que}}\\
&\deleted{\underset{\scriptscriptstyle\sum\alpha_i}{\underset{\scriptscriptstyle\Vert}{\zeta(\alpha)}}=2(\lambda-r)-r=2\lambda-3r}
\end{aligned}
\]
% archive sheet number: 50
